\documentclass[11pt,a4paper]{article}

% ---------------------------------------------------------------------------
% Supplementary Information for main.tex: the algorithm in full.
% Build: make   (see README.md)
% ---------------------------------------------------------------------------

\input{preamble}
\externaldocument[M-]{main}[main.pdf]

\renewcommand{\thesection}{S\arabic{section}}
\renewcommand{\thefigure}{S\arabic{figure}}
\renewcommand{\thetable}{S\arabic{table}}
\renewcommand{\theequation}{S\arabic{equation}}
\renewcommand{\thealgorithm}{S\arabic{algorithm}}
\renewcommand{\thelemma}{S\arabic{lemma}}
\renewcommand{\thetheorem}{S\arabic{theorem}}

\title{Supplementary Information\\[0.4em]
\large \swig: exact joint placement of V, D and J segments
for immune receptor germline assignment}

\author{Mateusz Kaduk}

\date{Draft for swig 0.39.7 --- \today}

\begin{document}
\maketitle

\noindent This supplement describes the algorithm of \swig{} 0.39.7 in
full: the model and the derivation of every prior term
(\cref{sec:model}), the exact recursion and its proof, per-read calls and
posteriors (\cref{sec:dag}), the plausibility constraints
(\cref{sec:guards}), the search with the guarantee of each step
(\cref{sec:search}) and the implementation (\cref{sec:implementation}).
\Cref{sec:protocol} gives the evaluation protocol: the data sets, the
scoring rules, the tool settings, the statistics and the timing conditions.
\Cref{sec:supresults} gives the results on the model's own terms that the
main text summarises, and \cref{sec:limitations} expands on the limitations.
Numbers prefixed with S refer to this supplement, unprefixed numbers to the
main text.

\tableofcontents

% ===========================================================================
\section{Model}
\label{sec:model}

\subsection{Problem}

Given a read $q$ and germline databases of V, D, J and optionally C alleles,
grouped into genes, \swig{} reports the strand of $q$, the alleles of its V,
D (zero or one) and J segments with their read and germline coordinates,
hence the N/P regions and the junction, and optionally the constant region.
The read may carry substitutions and short indels from SHM, sequencing
errors, trimmed segment ends, non-templated junction bases and unrelated
flanking sequence such as primers, adapters or UMIs.

\subsection{Segments are local alignments scored as log-odds}
\label{sec:scores}

A \emph{segment hypothesis}, an allele $a$ and a read interval $I$, is scored
by affine-gap local alignment~\cite{smith1981,gotoh1982} of $q[I]$ to $a$.
With match score $m$ and mismatch score $\mu$ a substitution column scores
$s(x,y)=m$ if $x=y$, $\mu$ if $x\neq y$, and $0$ if either base is ambiguous.
These are log-likelihood ratios~\cite{altschul1991,karlin1990}: let $\lambda$
solve
\begin{equation}
  \tfrac14 e^{\lambda m} + \tfrac34 e^{\lambda\mu} = 1 ,
  \label{eq:lambda}
\end{equation}
the Karlin--Altschul $\lambda$ for uniform base composition. Then
$e^{\lambda s(x,y)}$ is the likelihood ratio of ``$x$ was templated by $y$ at
divergence $\delta$'' against ``$x$ is random'', with
$1-\delta=\tfrac14 e^{\lambda m}$. For the default $(m,\mu)=(2,-3)$,
$\lambda=0.634$ nats per score unit and $\delta=11.2\,\%$: the familiar
$+2/-3$ nucleotide scores are the DNA analogue of a BLOSUM
matrix~\cite{henikoff1992} at that divergence. Gap runs of length $L$ cost
$g_o+(L-1)g_e$, with $(g_o,g_e)=(-9,-1)$ for V, $(-13,-1)$ for D and $(-17,-2)$
for J. N/P and flank bases score 0, since they are random under both
hypotheses.

V and J alignments are extended co-linearly, without gaps, towards the
germline's outer end (V~5$'$, J~3$'$) or the read's end, whichever comes
first. The extension stops before the column at which its running score
falls more than 30 below the alignment's edge (an
X-drop~\cite{altschul1997gapped} of 10 net mismatches), so an adapter or a
random 5$'$ flank is cut after about 17 bases on average (after 10 when
every base mismatches, later when chance matches occur) instead of being
taken whole for V.
Each stage keeps its best alignments both by extended and by local score, so
that a hit at the read's end, which pays for no flank, cannot take every slot
from a locally stronger hit that pays for a mutated one.

\subsection{Placements, cuts and paths}
\label{sec:paths}

A \emph{placement} is an aligned segment hypothesis. A \emph{path} is a
sequence of placements allowed by the grammar
\begin{equation}
  \mathrm{V} \;\to\; \mathrm{D}^{\{0,1\}} \;\to\; \mathrm{J} ,
\end{equation}
with partial paths V, J, VJ, VD and DJ allowed (never a lone D). Between two
consecutive placements there are two \emph{cuts}: the first keeps the read
bases before its \emph{exit} $x$, the second keeps the bases from its
\emph{entry} $y\geq x$ on, and the $y-x$ bases between them are N/P. Each
placement is \emph{clipped} to the alignment columns it keeps. Overlapping
placements may therefore split an overlap at any point, or leave part of it
to neither. Only the cuts order segments of different kinds: a D may precede
a J whose alignment starts before the D's, as long as the J is entered after
the D's exit (main-text \cref{M-fig:dag}). This is precisely the configuration in which a
sequential annotator lets a J alignment ``eat'' the D.

\subsection{Objective}
\label{sec:objective}

Read positions run from $0$ to $|q|$, and read intervals are half-open. A
placement $p$ has a kind $k(p)\in\{\mathrm V,\mathrm D,\mathrm J\}$, an
allele $a(p)$ of length $|a(p)|$ and a read interval
$[\operatorname{start}(p),\operatorname{end}(p))$. A \emph{match column} of
its alignment pairs one read base with one germline base, equal or not. A
read position $y$ is a \emph{canonical entry} of $p$ if a match column of $p$
holds read base $y$, and $x$ a \emph{canonical exit} if one holds read base
$x-1$; the \emph{clip} of $p$ to $[y,x)$ keeps the alignment columns from the
first of these two match columns to the second. A path is a sequence
\begin{equation*}
  \pi=\bigl((p_1,y_1,x_1),\dots,(p_n,y_n,x_n)\bigr)
\end{equation*}
of placements $p_i$, each with a canonical entry $y_i$ and a canonical exit
$x_i$, such that
\begin{enumerate}
  \item the word $k(p_1)\cdots k(p_n)$ is V, J, VJ, VD, DJ or VDJ;
  \item $y_1=\operatorname{start}(p_1)$ and $x_n=\operatorname{end}(p_n)$:
    the first placement is entered at its start and the last left at its
    end (V and J alignments are extended outward on purpose,
    \cref{sec:scores});
  \item $x_i-y_i\ge\ell_{\min}(p_i)$, the fewest read bases a clip of $p_i$
    may keep;
  \item $x_i\le y_{i+1}$ for $i<n$, the $y_{i+1}-x_i$ read bases between the
    two being N/P;
  \item with the open-flank cost on, a V immediately before a D and a J
    immediately after one each have a clip that keeps their outer end and is
    worth more than their presence cost (\cref{sec:dag}).
\end{enumerate}
Its value, in nats, is
\begin{equation}
\begin{aligned}
  v(\pi) \;&=\; \sum_{i=1}^{n} w(p_i,y_i,x_i)
     \;-\; c_{\text{start}}(p_1) \;-\; c_{\text{end}}(p_n),\\
  w(p,y,x) \;&=\; W_p(y,x) \;-\; C_{k(p)}
     \;+\; \ind[k(p)=\mathrm D]\,\bigl(\nu\,(x-y) - \tau_5\,t_5(p,y) - \tau_3\,t_3(p,x) + \kappa\bigr),\\
  c_{\text{start}}(p) \;&=\; \ind[k(p)=\mathrm D]\;\copen{V}\bigl(\operatorname{start}(p)\bigr),
  \qquad
  c_{\text{end}}(p) \;=\; \ind[k(p)=\mathrm D]\;\copen{J}\bigl(|q|-\operatorname{end}(p)\bigr),
\end{aligned}
\label{eq:objective}
\end{equation}
where $\ind[\cdot]$ is 1 when its condition holds and 0 otherwise, and
\begin{itemize}
  \item $W_p(y,x)$ is the worth of the columns of the clip of $p$ to
    $[y,x)$: $\lambda\,S_p(y,x)$ on the fixed scale, with $S_p(y,x)$ their
    alignment score (\cref{sec:scores}), and by default their worth at the
    read's own divergence (\cref{sec:perread});
  \item $C_k$ is the presence cost of a segment of kind $k$, derived from
    $P^{*}_k$ score units (\cref{eq:pd});
  \item for a D, $x-y$ is the number of read bases its clip keeps,
    $t_5(p,y)$ and $t_3(p,x)$ are the germline bases the clip leaves out at
    the 5$'$ and 3$'$ ends of $a(p)$, and $\nu$, $\tau_5$, $\tau_3$ and
    $\kappa$ are the N/P credit per kept base, the trim costs per left-out
    base and the credit of the D priors;
  \item $\copen{V}(\ell)$ and $\copen{J}(\ell)$ are the open-flank costs of
    a D with $\ell$ read bases before it, where its V is missing, or after
    it, where its J is missing (\cref{eq:openflank}).
\end{itemize}
No path is a lone D, so at most one of $c_{\text{start}}(p_1)$ and
$c_{\text{end}}(p_n)$ is non-zero. $W_p$ is a log-likelihood ratio and every
other term a log prior odds or log-likelihood ratio with a stated derivation
below; the call is the path of largest $v(\pi)$, under the tie order at the
end of this section. Of the constants the derivations leave free, the number
of junction positions $L=25$ and the even prior odds of a segment are chosen
rather than derived, the open-flank allowance and the credit $\kappa$ were
set on simulated reads, and the trim and N/P means are generic rather than
fitted (\cref{M-sec:discussion}).

\paragraph{Segment presence.} A D hypothesis is one of $N$ D genes placed at
one of $L$ junction positions, so its prior is spread over $NL$ alternatives.
A maximum a posteriori decision includes the best D, of alignment score
$S_{\mathrm D}$, when $\lambda S_{\mathrm D} > \ln(NL) - \ln\omega_{\mathrm D}$,
with $\omega_{\mathrm D}$ the prior odds of a D against none, which gives the
threshold, in score units,
\begin{equation}
  \pD \;=\; \frac{\ln(NL) - \ln \omega_{\mathrm D}}{\lambda}.
  \label{eq:pd}
\end{equation}
Genes are counted after identical sequences are merged. For the IMGT human
IGH set of \cref{sec:protocol-igsim} ($N=30$, $L=25$, even odds)
$\pD=10.45$: an exact
5-mer (score 10) loses and an exact 6-mer (12) wins. For TRB ($N=2$)
$\pD=6.17$. V and J pay the same kind of cost over their own genes, with the
same $L$ and even odds ($P^{*}_{\mathrm J}=7.91$ for the 6 IGHJ genes of that
set, $P^{*}_{\mathrm V}=12.54$ for its 113 IGHV genes). The presence cost
$C_k$ of \cref{eq:objective} is $P^{*}_k$ on the call's scale:
$\lambda\lfloor P^{*}_k\rfloor$ nats on the fixed scale and, by default,
$P^{*}_k/m$ matched columns valued at the read's divergence
(\cref{sec:perread}).
Without $C_{\mathrm J}$ and
$C_{\mathrm V}$, any positive clip completes or starts a path for free, so a read that
stops before its J or starts after its V receives a J or V from a few chance
bases. $P^{*}_k$ is kept unrounded; on the integer fixed scale charging
$\lfloor \pD\rfloor$ decides $S_{\mathrm D}>\pD$ exactly, since
$S_{\mathrm D}$ is an integer. A database of several loci
gives each locus its own costs. An alternative derivation counts the
hypotheses a local D search actually tries, every allele offset and read
position, through the Karlin--Altschul expectation $K w n_D e^{-\lambda S}$
($w$ read positions, $n_D$ allele offsets). It is available as an option
(\code{-{}-d-presence-rate}) but needs a D-presence rate that does not follow
from the search: at the D-presence rate of the IgSim \emph{mid} reads
(94.74\,\% carry a D remnant) $\pD$ falls to 9.86, below an exact 5-mer's
score of 10, and exact 5-mers win.

\paragraph{D trims and N/P lengths.} Let the clip of a D placement $p$ to
$[y,x)$ keep germline positions $[r_{\text{lo}},r_{\text{hi}})$ of its
allele $a(p)$: its match column at read base $y$ holds germline position
$r_{\text{lo}}$, and that at read base $x-1$ position $r_{\text{hi}}-1$. It
leaves out $t_5(p,y)=r_{\text{lo}}$ germline bases at the 5$'$ end and
$t_3(p,x)=|a(p)|-r_{\text{hi}}$ at the 3$'$ end. Under geometric trims of
means $\bar t_5=5$ and $\bar t_3=7$ each left-out base costs
$\tau=\ln((\bar t+1)/\bar t)$, i.e.\ $\tau_5=\ln(6/5)=0.18$ and
$\tau_3=\ln(8/7)=0.13$ nats. Under geometric N/P lengths of mean 8 a junction
pays $\ln(9/8)$ per N/P base, so every read base a D keeps is one N/P base
fewer and earns $\nu=\ln(9/8)=0.12$ nats. The D cost is credited with
$\kappa=\bar t_5\tau_5+\bar t_3\tau_3=1.85$ nats, the trim cost of a D with
the prior's mean trims. The trims then re-rank Ds while a typical D is called
as before, and a short exact match deep inside a long allele, which used to
be worth as much as an untrimmed short allele, pays for the germline it
leaves out. Each D term depends on one cut ($t_5$ and $\nu y$ on the entry,
$t_3$ and $\nu x$ on the exit), so $w(p,y,x)$ remains a difference of an
exit and an entry prefix value (\cref{eq:prefix}) and the exact recursion of
\cref{sec:dag} is unchanged.

\paragraph{Open flank.} A path whose last segment is a D asserts a junction
whose J the read does not show in the $\ell=|q|-\operatorname{end}(p_n)$
bases after the D; one whose first segment is a D asserts a V missed in the
$\ell=\operatorname{start}(p_1)$ bases before it. That D pays the
log-likelihood ratio of those bases under ``no junction, or a junction whose
J (V) was missed'' against ``a junction whose J (V) was missed''. For the
missing segment $S\in\{\mathrm V,\mathrm J\}$,
\begin{equation}
\begin{gathered}
  c^{S}_{\text{open}}(\ell) = \ln\!\Bigl(1 + \rho\,\frac{1-\varepsilon_S(h)}{\varepsilon_S(h)}\Bigr),
  \qquad h=\min\bigl(\max(\ell-N_{\text{open}},\,0),\,64\bigr),\\
  \varepsilon_S(h) = \max\Bigl(\varepsilon_{\min},\ \min_{0\le h'\le h}
     \Pr\bigl[X_{h'} \geq k_S(h')\bigr]\Bigr),\\
  k_S(h)=\Bigl\lfloor\frac{m_S\,h-A_S}{m_S-\mu_S}\Bigr\rfloor+1,
  \qquad X_{h} \sim \operatorname{BetaBinomial}(h,\,\alpha,\,\beta),
\end{gathered}
\label{eq:openflank}
\end{equation}
with $\Pr[X_{h'}\ge k_S(h')]=1$ whenever $m_S h'<A_S$, $h'=0$ included.
Here $h$ is the number of visible bases of the missing segment beyond the
free allowance $N_{\text{open}}$; $(m_S,\mu_S)$ are the segment's match and
mismatch scores and $A_S$ the smallest score its search stage keeps in a
64-base window, so that $h$ aligned bases with $k$ mismatches, scoring
$m_S h-(m_S-\mu_S)k$, stay below $A_S$ exactly when $k\ge k_S(h)$, and no
$h$ bases reach $A_S$ while $m_S h<A_S$; and $\operatorname{Beta}(\alpha,\beta)$
is the read's mismatch rate from its anchored V and J columns
($\alpha=k_a+\tfrac12$, $\beta=n_a-k_a+\tfrac12$ for $k_a$ mismatches among
$n_a$ columns). A present segment shows in $h$ bases if any prefix of them
shows it, so $\varepsilon_S(h)$ is an upper bound on the probability that a
segment present in the read stayed below the detection threshold. The cost
is 0 where $\varepsilon_S(h)=1$, as at $h=0$, non-decreasing in $\ell$ and
constant from $\ell=N_{\text{open}}+64$ on. The free allowance $N_{\text{open}}=24$ is
the largest 99.9th percentile (22--24 across simulated sets) of the distance
from a called D to the true J start, or from the V end to the D;
$\rho=\tfrac12$ (even odds) and $\varepsilon_{\min}=10^{-3}$ (the anchor
E-value gate) bound the cost to $[0,\,6.2]$ nats. At high divergence a
missing segment could well have been missed and the cost falls accordingly:
with 13 mismatches among 45 anchored columns (about 30\,\% SHM), $h=64$
visible bases cost 0.99 nats for a missing J and 0.71 for a missing V,
against 6.2 on a clean read. A real D lies within its N/P region of the read end, a chance D
anywhere in the search window, so the cost removes far chance Ds while
keeping real Ds of reads that stop at their D or start in CDR3
(\cref{sec:truncated}).

\paragraph{Order of the decision.} The adjudicator returns the path that
maximises, lexicographically, (1) its value $v(\pi)$ (\cref{eq:objective}),
by default at the read's own divergence (\cref{sec:perread}); (2) fewer D
segments, so a path with a D must
strictly beat the best path without one; and (3) germline \emph{support}, the
sum of its placements' unclipped scores. Key (3) separates paths that explain
the read equally well, typically two alleles of different 3$'$ length at a
junction base that can belong to either side; the allele with more support
admits more cuts with the same score and hence has the larger likelihood
summed over cuts. Co-optimal alleles, equal in clipped score and in matches
and mismatches on the called interval, are all reported. J names are
positional: a tied J allele is named only when it explains the called read
bases at the first allele's own germline coordinates (the same germline
interval), so IGHJ4*02 and IGHJ5*02 trimmed to their
shared 3$'$ part, which tie from different germline offsets, are not both
named; the other remains a tied alternative counted by the posterior.

\subsection{Significance}
\label{sec:significance}

Each V and J hit receives a Karlin--Altschul E-value~\cite{karlin1990}
$E=K w n e^{-\lambda S}$ against its database, with $w$ the length of the
searched read window and $n$ the database length, using exact ungapped
$\lambda$ and $K$ for the scoring in use. A hit is kept only at $E\le 1$, and a
read is \emph{called} only when its path contains an \emph{anchor}: a V or J
hit at $E\le 10^{-3}$ whose clip on that path still reaches the anchor score.
Chance calls on random sequence are therefore rare, of order $10^{-3}$ per
V or J search under the ungapped statistics; none of 2,000 random reads was
called (\cref{M-sec:specificity}). A J
downstream of a significant V needs only $E\le 1$, since given the V a J is
expected there.

% ===========================================================================
\section{Exact adjudication}
\label{sec:dag}

\subsection{Clipping and separability}

\begin{lemma}[clipping never helps]
\label{lem:clip}
Let $A$ be an optimal local alignment and $A[a..b]$ a contiguous
sub-alignment that begins and ends in match columns. Then
$S(A[a..b])\le S(A)$.
\end{lemma}
\begin{proof}
Every prefix and suffix of an optimal local alignment has non-negative
score, or removing it would give a higher-scoring local alignment. Hence
$S(A[a..b]) = S(A) - S(\text{prefix}) - S(\text{suffix}) \le S(A)$.
\end{proof}

Flatten a placement $p$ into its alignment columns $0,1,\dots$ and let
$\Pi_p[c]$ be the worth of its first $c$ columns. For a canonical entry $y$
let $c^{-}_p(y)$ be the index of the match column that holds read base $y$,
and for a canonical exit $x$ let $c^{+}_p(x)$ be one more than the index of
the match column that holds read base $x-1$. The clip to $[y,x)$ keeps
columns $c^{-}_p(y),\dots,c^{+}_p(x)-1$, so
$W_p(y,x)=\Pi_p[c^{+}_p(x)]-\Pi_p[c^{-}_p(y)]$: two table lookups and one
subtraction, exact because both end columns are match columns and a clip
therefore never splits a gap run. With the prefix values
\begin{equation}
\begin{aligned}
  \operatorname{entry}_p(y) &= \Pi_p[c^{-}_p(y)] + \ind[k(p)=\mathrm D]\,\bigl(\nu\,y + \tau_5\,t_5(p,y)\bigr),\\
  \operatorname{exit}_p(x) &= \Pi_p[c^{+}_p(x)] + \ind[k(p)=\mathrm D]\,\bigl(\nu\,x - \tau_3\,t_3(p,x) + \kappa\bigr),
\end{aligned}
\label{eq:prefix}
\end{equation}
the segment term of \cref{eq:objective} is
$w(p,y,x)=\operatorname{exit}_p(x)-\operatorname{entry}_p(y)-C_{k(p)}$, one
term per cut (\emph{separability}). $C_{k(p)}$, $c_{\text{start}}(p)$ and
$c_{\text{end}}(p)$ depend on no cut, and neither does the support key.

\subsection{The recursion}

Placements are the nodes of a lattice, sorted by segment kind (V, D, J) and
then by read start. Every transition between kinds goes forward in this
order wherever the two alignments start, and the cuts alone decide whether
it is feasible ($x\le y$). The paths through a node with a given D count and
first segment kind form a \emph{slot}, which holds one value per canonical
entry cut $y$ of the node: the best key of a path prefix entering there. By
condition~2 of \cref{sec:objective} a path can start at a node only at the
node's own start, in the slot whose first kind is the node's. The recursion
also carries two terms that are 0 by default and therefore left out of
\cref{eq:objective}: a linear N/P score of $\iota$ per N/P base and a pair
bonus $B_{\mathrm{VJ}}$ for a path that starts with a V and ends with a J,
both log prior odds set in score units and written here in nats ($\lambda$
nats per score unit, \cref{sec:perread}). With them it maximises
\begin{equation}
  v(\pi) \;+\; \iota\sum_{i=1}^{n-1}\,(y_{i+1}-x_i)
  \;+\; B_{\mathrm{VJ}}\,\ind\bigl[k(p_1)=\mathrm V,\ k(p_n)=\mathrm J\bigr],
  \label{eq:objective-ext}
\end{equation}
which is $v(\pi)$ at $\iota=B_{\mathrm{VJ}}=0$. The recursion is
\begin{align}
  X(b,x) &= \operatorname{exit}_b(x) - C_{k(b)} + \max_{y\le x-\ell_{\min}(b)}\bigl[F(b,y)-\operatorname{entry}_b(y)\bigr], \label{eq:rec-exit}\\
  F(e,y') &= \max\Bigl(\max_{b\to e}\,\max_{x\le y'} \bigl[X(b,x) + \iota\,(y'-x)\bigr],\; s_e(y')\Bigr), \label{eq:rec-entry}\\
  \text{best} &= \max_{(g,b)}\,\bigl[X\bigl(b,\operatorname{end}(b)\bigr) + B(g,b) - c_{\text{end}}(b)\bigr], \label{eq:rec-best}
\end{align}
where $y$ ranges over the canonical entries of $b$, $x$ over its canonical
exits and $y'$ over the canonical entries of $e$; $F$ and $X$ are kept per
slot, which the notation leaves implicit, and a maximum over no candidates is
$-\infty$. In \cref{eq:rec-entry}, $b\to e$ ranges over the predecessors
whose slot and kind may precede $e$ (the next kind continues a word of
condition~1, condition~5 holds, and $e$'s slot has $b$'s first kind and $b$'s
D count plus $\ind[k(e)=\mathrm D]$), and $s_e(y')=-c_{\text{start}}(e)$ if
$y'=\operatorname{start}(e)$ in the slot in which a path starts at $e$, and
$-\infty$ otherwise. In \cref{eq:rec-best}, $(g,b)$ ranges over the nodes $b$
and the slots $g$ of $b$ in which a path may end at $b$ (the kinds form a
word of condition~1), and $B(g,b)=B_{\mathrm{VJ}}$ if $g$'s first kind is V
and $b$ is a J, 0 otherwise. Because each term depends on at most one cut,
and the N/P term $\iota(y'-x)$ splits into $\iota y'-\iota x$:
\begin{enumerate}
  \item the best entry for an exit, $\max_{y\le x-\ell_{\min}(b)}[F(b,y)-\operatorname{entry}_b(y)]$,
    is a running maximum along the node;
  \item the best exit for an entry, $\max_{x\le y'}[X(b,x)-\iota x]$, is a
    second running maximum along $b$, built once per slot of $b$ and per
    successor kind, not per edge;
  \item every entry $y'\ge\operatorname{end}(b)$ sees the same best exit of
    $b$, so the edge $b\to e$ adds it once at $e$'s first entry at or after
    $\operatorname{end}(b)$, and a prefix maximum over $e$'s entries carries
    it to the rest.
\end{enumerate}
An edge therefore costs $O(\text{overlap}+1)$ and a read
$O\bigl(\sum_{\text{slots}}\text{len}+\sum_{\text{edges}}(\text{overlap}+1)\bigr)$,
plus one scan of each node's predecessor slots; \cref{alg:dag} gives the
recursion as pseudocode.
With the open-flank cost on, condition~5 bars one transition when the
lattice opens its slots: a V or J none of whose clips that keep its outer end
is worth more than its presence cost may not immediately precede or follow a
D, since such a placement is no evidence of the junction's V or J (a junk J
of a base or two would otherwise close a D's flank for less than the cost).
The open-flank costs $c_{\text{start}}(e)$ and $c_{\text{end}}(b)$ are
constants of the node and the read, added once where a path starts and once
where it ends, and the D priors enter a D node's entry and exit prefix
values (\cref{eq:prefix}) once, when the node is laid out; neither breaks
separability. The recursion is a semi-Markov Viterbi~\cite{viterbi1967,yu2010hsmm}:
segments of variable length scored as wholes, with the cut positions as the
hidden boundaries.

\begin{theorem}
\label{thm:dag}
The recursion of \cref{eq:rec-exit,eq:rec-entry,eq:rec-best} returns the
maximum of \cref{eq:objective-ext}, and so at the default
$\iota=B_{\mathrm{VJ}}=0$ the maximum of $v(\pi)$ in \cref{eq:objective},
over all paths of \cref{sec:objective} (every placement, every grammatical
word and every pair of canonical cuts that satisfy its conditions 1--5),
with the tie order of \cref{sec:objective}.
\end{theorem}
\begin{proof}[Proof sketch]
Every grammatical predecessor of a node precedes it in the order by kind, so
by induction over that order $F(e,y)$ is the best value of a path prefix
that enters $e$ at $y$, and $X(b,x)$ that of a prefix that leaves $b$ at $x$.
A path's value is a sum of terms each of which depends on at most one cut
(\cref{eq:prefix}, with $\iota(y'-x)=\iota y'-\iota x$), so the maximum over
the pair $(x,y')$ of an edge factorises into the two running maxima above,
and the maximum over the entries beyond $b$'s end into one shifted prefix
maximum. Ties are resolved by comparing the lexicographic key
(value, fewer D, support), each of whose parts is separable in the same way.
\end{proof}

\begin{algorithm}[t]
\caption{Exact adjudication of one read. The posterior pass runs the same
loop in the sum-product semiring, with prefix sums in place of running
maxima.}
\label{alg:dag}
\small
\begin{algorithmic}[1]
\Require placements of the read sorted by kind (V, D, J), then read start;
  for each placement $b$ the prefix values $\operatorname{entry}_b(y)$ and
  $\operatorname{exit}_b(x)$ of \cref{eq:prefix} at its canonical cuts
\For{each node $e$ in order, and each slot of $e$ (D count, first kind)}
  \State $F(e,y)\gets-\infty$ for every entry $y$
  \If{this is $e$'s start slot and $e$ may start a path}
    \State $F(e,\operatorname{start}(e))\gets -c_{\text{start}}(e)$
  \EndIf
  \For{each predecessor $b$ whose slot and kind may precede $e$}
    \State $F(e,y)\gets\max\bigl(F(e,y),\,M_b(y)+\iota y\bigr)$ for every entry $y<\operatorname{end}(b)$
    \State offer $M_b(\operatorname{end}(b))+\iota y_0$ at the first entry $y_0\ge\operatorname{end}(b)$
  \EndFor
  \State carry the offers to later entries by a prefix maximum, $+\iota$ per base
  \State $G(x)\gets\max_{y\le x-\ell_{\min}(e)}\bigl[F(e,y)-\operatorname{entry}_e(y)\bigr]$ \Comment{running maximum along $e$}
  \State $X(e,x)\gets\operatorname{exit}_e(x)-C_{k(e)}+G(x)$ for every exit $x$
  \State $M_e(x)\gets\max_{x'\le x}\bigl[X(e,x')-\iota x'\bigr]$, once per kind that may follow $e$
\EndFor
\State \Return $\max_{(g,b)} X(b,\operatorname{end}(b))+B(g,b)-c_{\text{end}}(b)$ and its trace-back
\end{algorithmic}
\end{algorithm}

The search runs over every canonical entry of the successor, not only those
inside the overlap. Restricting entries to the overlap appears justified by
\cref{lem:clip}, but the lemma bounds prefixes that start at the alignment's
own start, whereas the bases between $b$'s end and a later entry are a
middle stretch of $e$'s alignment and can score below zero: two mismatches
right after the V's end, as a mutated read has, are better left to N/P.
Likewise, ordering nodes by read start alone excludes a D before a J whose
alignment starts earlier. Both restrictions lose optimal paths. \swig's test
suite therefore checks \cref{thm:dag} against a
memoised exhaustive enumeration of every grammatical path and canonical cut
pair, on synthetic reads, on the pipeline's own placements for hypermutated
reads, under a two-D grammar (VDDJ optima occur), with the per-read model at
five divergences, with the open-flank cost and with the D priors, and on
constructed reads that each restricted search gets wrong.

\subsection{Calls at the read's own divergence}
\label{sec:perread}

The scores $(2,-3)$ are log-odds only at $11.2\,\%$ divergence. A fixed scale
makes unmutated reads under-confident, since a single nucleotide difference
between two alleles is worth $\ln(3(1-p)/p)=8.0$ nats at $p=0.1\,\%$ rather
than $\lambda\cdot 5=3.2$, and heavily mutated reads over-confident. \swig{}
therefore estimates the read's divergence $\phat$ before adjudication from
the columns no path contests---the best V's match columns before every D and
J placement and the best J's after every V and D placement, N columns
excluded, or all columns of the two placements when fewer than 50
remain---as the Jeffreys predictive rate $(k+\tfrac12)/(n+1)$ for $k$
mismatches among $n$ columns~\cite{jeffreys1946}, clamped to
$[10^{-4},\,0.5]$. Apart from that fallback the junction bases are not among
these columns, so they are not used twice. A substitution column of score
$s\neq 0$ is then worth
$\lambda_r s+\mu_r$ nats, with $(\lambda_r,\mu_r)$ chosen so that a match is
worth $\ln 4(1-\phat)$ and a mismatch $\ln(4\phat/3)$; at $\phat=\delta$,
$\lambda_r=\lambda$ and $\mu_r=0$. Gaps, $\iota$ and the pair bonus stay at
the prior $\lambda$, since $\phat$ measures substitutions.

Every term of \cref{eq:objective} is then an integer number of units of
$1/F$ nats, $F=2^{16}$. In $W_p$, a substitution column of score $s\neq0$ of
a placement of kind $k\in\{\mathrm V,\mathrm D,\mathrm J\}$ is worth
$s\,u_k+v_k$ units, with $u_k=\operatorname{round}(F\lambda_r)$ and
$v_k=\operatorname{round}(F\mu_r)$ and $(\lambda_r,\mu_r)$ computed from that
kind's own match and mismatch scores (identical for the default $(2,-3)$); an
N column is worth 0, and a gap run of score $g$ is worth $g\,u_0$ units, with
$u_0=\operatorname{round}(F\lambda)$, as are $\iota$ and $B_{\mathrm{VJ}}$
per score unit. Each presence cost is read as $P^{*}_k/m$ matched columns of
its kind rather than as a constant number of nats, with $m$ the match score
(2 for every kind), and a matched D column also earns the N/P credit $\nu$:
\begin{equation}
\begin{gathered}
  C_{\mathrm V}(\phat)=\Bigl\lfloor P^{*}_{\mathrm V}\,\frac{m u_{\mathrm V}+v_{\mathrm V}}{m}\Bigr\rfloor,\qquad
  C_{\mathrm J}(\phat)=\Bigl\lfloor P^{*}_{\mathrm J}\,\frac{m u_{\mathrm J}+v_{\mathrm J}}{m}\Bigr\rfloor,\\
  C_{\mathrm D}(\phat)=\Bigl\lfloor \pD\,\frac{m u_{\mathrm D}+v_{\mathrm D}+\operatorname{round}(F\nu)}{m}\Bigr\rfloor
\end{gathered}
  \label{eq:presence-read}
\end{equation}
units. The D terms $\nu$, $\tau_5$, $\tau_3$ and $\kappa$ are charged as
$\operatorname{round}(F\nu)$, $\operatorname{round}(F\tau_5)$,
$\operatorname{round}(F\tau_3)$ and $\operatorname{round}(F\kappa)$ units,
and an open-flank cost $c$ as $\operatorname{round}(Fc)$. The maximum and
every tie are therefore exact, and the recursion and \cref{thm:dag} carry
over unchanged.

The fixed scale (\code{-{}-calls fixed}) instead takes
$W_p(y,x)=\lambda S_p(y,x)$, $C_k=\lambda\lfloor P^{*}_k\rfloor$ and
$\nu=\tau_5=\tau_3=\kappa=0$, and charges an open-flank cost $c$ as
$\lambda\Delta$ with $\Delta=\lfloor\pD+c/\lambda\rfloor-\lfloor\pD\rfloor$
score units, so that $v(\pi)/\lambda$ is an integer and an integer D score
$S_{\mathrm D}$ beats $\lfloor\pD\rfloor+\Delta$ exactly when it beats
$\pD+c/\lambda$. At $\phat=\delta$ every per-read column is its fixed score
times $u_0$, so with integer $P^{*}_k$ and the D priors and open-flank cost
off, the fixed scale is an exact special case of the per-read one.

The D cost is a threshold against
chance D hits in random N/P sequence, whose null distribution does not depend
on $\phat$ and is dominated by exact runs. Read in matched columns it keeps
the fixed scale's admitted classes of exact runs at every divergence for a D
with the prior's mean trims (runs of 5 lose, runs of 6 win when their trims
cost $\kappa$); the trim terms move runs that are trimmed less or more than
that across the threshold. A cost constant in nats would instead admit
exact 5-mers with the prior's mean trims on unmutated reads.

\subsection{Posterior probabilities}
\label{sec:posterior}

The objective is a log-likelihood ratio in nats, so it also defines a
distribution over the paths $\pi$ of \cref{sec:objective}, each with its
placements and canonical cuts,
\begin{equation}
  P(\pi\mid q) = \frac{\exp v(\pi)}{Z},\qquad Z=\sum_{\pi'}\exp v(\pi'),
\end{equation}
with $v$ of \cref{eq:objective} (\cref{eq:objective-ext} when $\iota$ or
$B_{\mathrm{VJ}}$ is set) on the call's scale.
Forward--backward~\cite{rabiner1989,durbin1998} over the same lattice---the
sum-product twin of \cref{eq:rec-exit,eq:rec-entry,eq:rec-best}, with prefix
and suffix sums where the Viterbi pass has running maxima---gives the marginal
of every placement and of every D count in the same $O(\text{overlap}+1)$ per
edge. Under the default call both passes read the same integer model, so the
Viterbi path is a mode of the reported distribution. A segment's reported
posterior is the mass of every placement of its kind that is tied with the
call on the called interval or explains most of the called bases with
germline the primary allele contains; the event is ``the reported names
include the read's allele, up to germline equivalence'', which is exactly
what the evaluation scores (\cref{sec:protocol-igsim}). The read's D-presence
probability is $1-P(\text{no D})$.

The arithmetic is linear with one log-scale per slot, as in scaled HMM
recursions. Cut weights are running products along the alignment, so no
transcendental function is evaluated per cut; entry and exit weights are
centred separately; the N/P factor is accumulated as a relative recursion
whose factors are $\le 1$; and the cost of a gap between non-overlapping
neighbours joins the slot's log-scale rather than multiplying it. A read
whose evidence would leave the range of double precision (a placement
spanning more than 1,200 nats on one side, or an overlap gaining more than
600) receives no posterior rather than a wrong one. Tests compare the
marginals with exhaustive summation under VJ, VDJ and VDDJ grammars, with
N/P costs, N columns, insertions, long costly gaps, the open-flank cost and
the D priors.

An optional decoder (off by default) takes a call with exactly one D and
finds the allele of largest posterior marginal,
$a^{*}=\argmax_a\sum_{p\ni a}P(p)$ over D placements $p$; the marginal is
not conditioned on the call's topology. When $a^{*}$ beats the call's
allele, the path is re-traced with the call's first and last kind and D
count on the lattice masked to the V and J placements and $a^{*}$'s D
placements, and kept when it stays anchored and its reported D posterior is
no lower.


% ===========================================================================
\section{Plausibility constraints}
\label{sec:guards}

On full-length simulated reads the optimum of \cref{eq:objective} almost
always pairs a V and a J that could form one rearrangement, but on
truncated, chimeric and cross-locus reads it need not. \swig{} applies, as
constraints outside the objective and on by default, the plausibility checks
of SwiftIG, an earlier C++ aligner (unpublished) whose method note \swig{}
began as a Rust implementation of: minimum hit lengths of 24 (V) and 10 (J)
local alignment columns; a V--J pairing rule (the V and J align both
conserved codons in order, or hold at least 80 local alignment columns
together and either face their junction or hold a long split), to which
\swig{} adds that, unless the pair aligns both conserved codons in order, a
V shorter than 80 columns pairs only with a J that is significant on its
own; and a span of at most 225 bases between the V's
exit and the J's entry. SwiftIG called a D only between a V and a J that
pair. \swig{} also admits a D beside a V or a J alone, but only when it pays
the open-flank cost and lies within 60 bases of the segment the read shows.
The call is the unconstrained optimum when it satisfies the constraints;
otherwise the lattice is re-adjudicated over subsets of placements (the V
placements alone, the J placements alone, and each J with the V placements
it pairs with, with and without the Ds; under the open-flank cost also V and
D placements within reach of each other, and D and J placements), and the
call is the best admissible subset optimum. Within each subset the key is the
exhaustive optimum, but a subset whose optimum breaks a constraint
contributes nothing, so the guarded call is \emph{defined} rather than exact.
On the test suite's 435 checked calls, no admissible path over any choice of
at most one V, D and J among a read's best 4, 8 and 4 placements beat the
call under the per-read model (one did under the fixed scale). Junctions are
reported only from conserved anchors
(Cys104, Trp/Phe118) that the called alignments actually align, with a
bounded recovery of a Cys104 lying just past a clipped V end.

% ===========================================================================
\section{Search}
\label{sec:search}

Each stage narrows the hypotheses the lattice sees. \Cref{tab:search} states
what each step can lose: \emph{exact}; \emph{lossless}, which provably drops
nothing the model could choose; \emph{bounded}, which drops only what lies
outside a stated geometric bound; \emph{defined}, a stated selection rule
whose output is the definition of the result rather than an approximation
of an optimum; \emph{approximate}, a known gap to the model's optimum; and
\emph{heuristic}, which is measured, not proven. \Cref{M-fig:pipeline} of the main text
shows the flow.

\begin{table}[tp]
\centering
\caption{Search steps and what each can lose.}
\label{tab:search}
\footnotesize
\begin{tabularx}{\linewidth}{@{}l l X@{}}
\toprule
Step & Guarantee & Rule and justification \\
\midrule
Strand & heuristic & Count read positions whose 11-mer occurs in the V/J germline on each strand; trusted when the winner exceeds its Poisson null $\lambda_0+3\sqrt{\lambda_0}$ ($\lambda_0$: the read's $k$-mer count times the germline's 11-mer density), otherwise compare V seed votes on both strands. \\
V retrieval & heuristic & Seed votes per gene frame (9-mers weight 4; 7-mers weight 1 while the leader holds fewer than 25 nine-mer seeds). Keep 3 genes, and up to 16 while their votes are within $z\sqrt{2\cdot 4\cdot L}$ of the leader's $L$, $z=2$: about two standard deviations of the difference of two seed counts, so the margin is proportionally wider for weakly supported, mutated reads. Every allele of a kept gene is a candidate. \\
V band & bounded & The seed cluster's diagonal range $\pm 8$. A frame's runner-up cluster is taken as the far side of an indel, and the band spans both, when the two span at most 64 diagonals, the runner-up holds at least the votes of 3 nine-mers and 1/8 of the frame's best cluster, and that cluster those of 25. Holds unless a longer indel is not joined this way: its far side is unseeded or below those shares (e.g.\ an indel 20 bases from a V end), the two span more than 64 diagonals, or the frame's best cluster is weakly seeded. \\
V, J scoring & lossless & One representative $R$ per gene frame is aligned; a sibling $A$ satisfies $\underline{L}\le S^{*}(A)\le S^{*}(R)+\Delta(R,A)$ and is skipped when the upper bound is below the third-largest lower bound. The ranking from the third score up equals full scoring. \\
V/J keep & defined & Top 3 by extended and top 3 by local score with $E\le1$, plus ties. \\
J & bounded & Every J allele of the anchoring V's locus, unbanded, from 12 bases before the earliest V end to at most 225 bases plus the longest J past the last anchoring V's end; the whole read when no V is an anchor. \\
D windows & bounded & $[\text{V end}-12,\ \text{J start}+12)$ from the best anchoring V and the best J placement, plus the conserved-anchor window $[\text{Cys104}+3,\ \text{Trp/Phe118})$; an open side reaches 60 bases. \\
D prefilter & lossless & An allele is aligned only if it shares an $\ell$-mer with the window, $\ell$ derived from the admission gate (exact run $\ge 4$). \\
D alignment & approximate & One optimal local alignment per allele and window, the best 32 admissible ones (plus ties) per window; its clip at the chosen cuts need not be the best clip of that allele. \\
Adjudication & exact & \Cref{thm:dag}. \\
Guards & defined & Best admissible subset optimum (\cref{sec:guards}). \\
C & exact within 7 bases & Anchored DP at the J exit (\cref{sec:constant}). \\
\bottomrule
\end{tabularx}
\end{table}

\paragraph{Gene frames and the index.} The alleles of a gene share almost
all seeds. An allele joins its gene's coordinate frame when at least 30\,\%
of its seeds occur in the gene's longest allele and every one of those
shared seeds agrees on one offset; SNP and truncated alleles join, an allele with an
internal indel gets its own frame. The $k$-mer index is built over frames in
compressed-sparse-row form, with 2-bit packed keys used directly as array
indices, so a seed shared by twenty alleles is one posting and one vote and
retrieval work scales with genes rather than alleles. Votes are packed into
64-bit integers and radix-sorted in linear time; one scan then yields the
best diagonal cluster per frame. Seeds with more than 256 postings are
skipped, and a seed repeated within a frame does not vote, so repeats cannot
collect quadratic votes. We measured closed and open syncmers and
minimisers~\cite{edgar2021syncmers,shaw2022syncmers} as seed samplers
(minimisers only under an earlier fixed gene margin). Under the noise-based
margin every $k$-mer beat closed syncmers by 1.1 points of mean V gene
accuracy (4.4 at 30\,\% SHM) for 0.4\,\% more instructions, and beat open
syncmers too, so it is the default.

\paragraph{Sibling bounds.} For a sibling $A$ aligned in the same frame
diagonals as its representative $R$, $\Delta(R,A)$ sums the best profile gain
at each frame position both alleles cover and the best substitution score
where only $A$ does, so $S^{*}(A)\le S^{*}(R)+\Delta(R,A)$. The lower bound
$\underline{L}$ rescores a traced representative path clipped to $A$'s span
with $A$'s profile, in time proportional to the positions where the two
differ. A sibling whose bounds meet needs no alignment, and one whose lower
bound already reaches the cut-off is aligned outright with its traceback
reused. Because the cut-off is a lower bound on the true third score, the
kept ranking is that of aligning every allele.

\paragraph{D search.} The junction windows reclaim 12 bases into V and J,
and the conserved-anchor window depends on where the J's Trp/Phe118 lands,
not on how far its alignment reaches left; with the lattice's ordering by
kind, the read's own J cannot remove its D from consideration. A chance J
placement between the V and the D of a read that stops at its D can, since
the window ends 12 bases after the start of the best J placement, used or
not. Each D
allele sharing an $\ell$-mer with a window is aligned there, unbanded, and
its optimal local alignment competes for the 32 places (plus ties) that the
window's admissible D hits have in the lattice. D admission requires an
exact run of at least 4 bases; 3, 5 and 6 were measured and are worse.

\subsection{Constant region}
\label{sec:constant}

With a C database the constant region is assigned after adjudication, from
finished calls only, so V, D and J calls are unchanged by construction. J and
C are joined by an RNA splice at fixed phase, so C's germline base 0 lies at
the called J's exit $x$, up to reference conventions and indels near the J's
3$'$ end; hypothesis $H(a,o)$ places base 0 of allele $a$ at $x+o$,
$|o|\le 7$. An alignment anchored at $C[0]$ on one of these 15 diagonals is
computed exactly by any local aligner run on $[N]+C$ whose first column
scores a bonus larger than any alignment avoiding it, so the SIMD kernel
serves unchanged. Alleles whose sequences agree on the reachable prefix form
reach classes and are aligned once (the IMGT human IGHC set used here has 105
distinct alleles but only 13 distinct 32-nt prefixes). C is called when its
anchored E-value $E_C=15\,N_C\,e^{-\lambda S}\le 10^{-3}$, with $N_C$ the
number of C genes of the J's locus, from a Lundberg-type bound for the
walk from each start, on at least 12 aligned columns at 90\,\% identity, and
its posterior is computed at C's own divergence, since C is not
hypermutated.

\subsection{Germline annotation}

Conserved anchors fix every reading frame and junction end: V Cys104 and J
Trp/Phe118. They and the V region boundaries are settled once per database,
in order of trust: header keys, IMGT-gapped numbering when the gaps are
consistent, a joint motif search for Cys23, Trp41 and Cys104, transfer from
the most similar allele by local alignment, and finally database models: a
Jeffreys-smoothed model of the FR3 length given the Trp41--Cys104 spacing,
and of the Cys23 position. Every derived value that may be wrong carries an
\emph{uncertain} flag, so an error without the flag is a silent error. Each
allele's locus is read from its name or inferred from its relatives and,
failing that, from an embedded panel of human and mouse family
representatives; the locus decides which alleles may form one
rearrangement and which gene counts price a segment. V alleles named
without the \code{*} allele separator join the gene of their nearest named
relative ($\ge 0.95$ of the smaller self-score; the rest are grouped among
themselves), so the gene-level retrieval budget does not silently become an
allele cap.


% ===========================================================================
\section{Implementation}
\label{sec:implementation}

\swig{} is a Rust crate (edition 2024, Rust $\ge 1.89$) with no dependencies:
about 38,000 lines of library code with inline unit tests, 5,400 lines of
integration tests and 1,300 of examples and WebAssembly bindings. Every
pipeline component is a statically dispatched trait---aligner, retrievers,
candidate budget, D window policy, admission, adjudicator, grammar and path
model---so swapping one changes the type and the pipeline monomorphises
without trait objects. The substitution model is also a trait, applied once
when the database is prepared to build the per-allele profiles that the
kernels read as data.

\paragraph{Inter-sequence SIMD.} Following Rognes~\cite{rognes2011}, each
16-bit lane runs the scalar Gotoh recurrence for a different read--allele job
rather than striping one alignment~\cite{farrar2007}. The kernel is written
once over a lane abstraction and compiled for AVX-512BW (32 lanes), AVX2,
NEON, WebAssembly~\cite{haas2017wasm} \code{simd128} and portable lanes (16).
By default the fastest of AVX2, NEON and portable lanes is chosen at run time;
AVX-512BW runs only on request (\code{-{}-simd avx512}), since 512-bit
instructions lower the clock of many CPUs. On WebAssembly the backend is
fixed at build time (\code{simd128}, or the scalar kernel without it).
Banded jobs use band-relative coordinates, so lanes whose
bands start on different diagonals walk the same grid~\cite{chao1992band}.
When every lane's profile depends on base identity alone a cell's
substitution is one comparison and one blend; otherwise five profile rows are
selected by mask, exact for any substitution model. The horizontal gap state
uses $D'=\max(t+g_o,\,D+\max(g_o,g_e))$ with $t=\max(0,\text{diag},I)$, which
shortens the loop-carried dependency to one add and one maximum and is exact
for any gap costs. Arithmetic saturates in 16 bits; a job enters a lane only
if its attainable local score is at most 30,000, no step costs less than
$-1{,}000$ and its dimensions fit 16-bit positions, so saturation never clips
a real cell, and other jobs run on the scalar kernel.
Traceback codes are bit-identical to the scalar kernel's, so one traceback
walk serves both.

\paragraph{Batching and determinism.} Each stage runs over a chunk of 64
reads, giving one score batch and a few traceback waves per stage, while
every decision stays per read; a read's result does not depend on its chunk,
which is tested. Threads self-schedule chunks from an atomic counter with
one workspace each, so there is no allocation in the DP after warm-up, and
results return in input order. The scalar kernel equals a textbook Gotoh,
every SIMD backend equals the scalar kernel in scores, end cells and
tracebacks, the batched pipeline equals the per-read scalar pipeline, and
output is byte-identical across backends and between x86-64 and AArch64.

\paragraph{Input, output and interfaces.} Reads are FASTA, FASTQ or an AIRR
Rearrangement table, streamed in batches of 16,384; identical sequences are
assigned once per batch, and across batches through a cache of the 4,096
most recently assigned distinct sequences, which changes no output. Output is
the AIRR Rearrangement schema~\cite{vanderheiden2018airr} with the
posteriors, E-values, alternatives and region coordinates. The crate provides
a command line (\code{swig}, with \code{prepare}, \code{simulate} and
\code{evaluate} subcommands) and a Rust library; a companion workspace crate
(\code{swig\_wasm}) builds a WebAssembly module with a C ABI. All three share
one run implementation.

\paragraph{Tests.} The default suite of 350 tests, all passing on the
version described (one further long-running accuracy test, of 20,000 reads
per preset, runs on demand and is not counted among them), includes the exhaustive oracles for \cref{thm:dag} and the
posterior marginals, kernel equivalence on random batches for every backend
the host runs (AVX-512BW, AVX2 and portable lanes on x86-64; NEON and
\code{simd128} on AArch64 and WebAssembly builds), the guards on constructed
reads, robustness on
degenerate input (empty, all-N, IUPAC, low-complexity, 20\,kb and chimeric
reads, and every prefix and suffix of a read), and accuracy floors on
simulated reads. It runs on release code with debug assertions and overflow
checks enabled.

% ===========================================================================
\section{Evaluation protocol}
\label{sec:protocol}

This section gives the data sets, the scoring rules, the tool settings, the
statistics and the timing conditions behind the results of the main text.
All reads are simulated, so that every call can be scored against a known
answer. Two independent simulators are used, GenAIRR and IgSim, and \swig{}
0.39.7, IgBLAST 1.22 and MiXCR 4.7.0 are compared on the reads of both.

\subsection{GenAIRR reads and the tie-set score}
\label{sec:protocol-genairr}

We use reads simulated by GenAIRR
2.3.1~\cite{konstantinovsky2024genairr} in two groups of data sets.

\paragraph{Rhesus panels.} The first group comes from the IgBench harness
(\url{https://github.com/mashu/IgBench.jl}). We use its bundled
\code{rhesus\_igh} experiment, which simulates from rhesus macaque IGH
alleles. It consists of two panels of 50,000 reads, one with IgM-like and
one with IgG-like rates of SHM under the S5F model~\cite{yaari2013s5f}. Both
panels have a polymerase indel probability of 0.003 and a secreted constant
region placed on the simulated sequence.

\paragraph{Human sets.} For the second group we simulated 11 sets of 10,000
reads from GenAIRR's built-in human references (seed 7, forward strand). The
IGH reference, from OGRDB~\cite{lees2020ogrdb}, has 198 V, 33 D and 7 J
alleles; IGK has 168 V and 8 J alleles, IGL 181 V and 10 J, and TRB 98 V, 3
D and 16 J. There are five IGH sets: an unmutated set, three sets with S5F
mutation at 2, 5 and 10\,\%, and a \emph{realistic} set. The realistic set
combines S5F mutation at 5\,\% with the artefacts of library preparation and
sequencing: a 5$'$ primer trim of 0--15 bases, a 3$'$ trim of 0--5 bases,
0--1 polymerase indels, 0.2\,\% PCR errors, 0.2\,\% sequencing errors and
0--2 N calls. IGK and IGL each have a 5\,\% SHM set and a realistic set. TRB
has an unmutated set and a realistic set with the same artefacts but without
SHM.

The D trim and N/P-length priors were designed on the IGH sets of this grid
(\cref{sec:dpriors}). As an out-of-sample check, four of the IGH scenarios
were therefore simulated again from a disjoint part of GenAIRR's random
stream (seed 50000; \cref{tab:oos}). In both groups, every tool assigns
against the same V, D and J FASTA files exported from GenAIRR.

\paragraph{Truth and score.} A read does not always determine its allele:
after trimming and mutation, several alleles can match it equally well.
GenAIRR therefore records the truth of each segment as a \emph{tie set}.
Each allele earns a point for every read base, within the true region of the
segment in the read, that matches the allele at the same germline index as
in the sampled allele. The tie set contains every allele with the maximal
number of points. The call of a tool for a segment is likewise a set of
alleles, $P$. Against the tie set $G$ it scores the Jaccard index
$|P\cap G|/|P\cup G|$ per read and segment; two empty sets score 1 and a
single empty set scores 0. This is the allele accuracy of IgBench, and the
tables give its mean in per cent. One correct allele against a tie of two
alleles therefore scores $\tfrac12$, and so does a call of two alleles
against a truth of one.

\paragraph{Scoring MiXCR.} IgBLAST and \swig{} list all co-optimal alleles
of a segment. MiXCR's AIRR output names a single allele per segment, which
the Jaccard index cannot credit in full when the truth is a tie. MiXCR is
therefore scored in two ways. Its \emph{AIRR call} is the single allele of
its AIRR output; this is what IgBench itself scores. Its \emph{tied hits}
are all the alleles whose alignment score equals MiXCR's top score for the
segment, obtained by exporting every hit with its score
(\code{exportAlignments}). The tied hits give the fairer comparison with
IgBLAST and \swig, and the one that is more favourable to MiXCR.

\subsection{IgSim reads and germline-equivalent accuracy}
\label{sec:protocol-igsim}

\paragraph{Simulator and presets.} \swig{} includes a Rust port of the
generative model and presets of IgSim.jl~\cite{igsimjl}. It chooses alleles
uniformly, trims the segment ends as exonucleases do, adds N nucleotides and
random flanks, applies SHM and single-base indels, each to a
preset-dependent share of the reads, and adds instrument substitutions.
Every read base is labelled with its germline origin. Three presets set the
level of mutation. The \emph{easy} preset has no SHM. The \emph{mid} preset
is IgM-like: 40\,\% of the reads carry SHM, at rates of 0.5--7\,\%, and
1-nt indels are switched on for 8\,\% of the reads at up to 0.2\,\% per
base, so that about 2\,\% of the reads carry one. The \emph{hard} preset is
IgG-like: 97\,\% of the reads carry SHM, with a median rate of about 7\,\%,
and indels are switched on for 15\,\% of the reads at up to 0.4\,\% per
base, so that about 7\,\% carry one. To probe the extreme, the \emph{hard}
preset is also run with the SHM rate fixed at 30\,\% on every read.

\paragraph{Germline.} The germline is the human IGH set of
IMGT/GENE-DB~\cite{giudicelli2005genedb}, release 202632-7. It has 480 V, 48
D and 15 J alleles, which \swig{} reduces to 457, 41 and 15 by merging
identical sequences. Every tool receives the set as plain FASTA; \swig{}
annotates the conserved anchors itself (\cref{sec:search}). The simulator
draws only V and J alleles with a trusted anchor, since a junction truth
built on a guessed anchor would rest on the annotation that it is meant to
test.

\paragraph{Correct calls.} An allele call is correct when it is
\emph{germline-equivalent} to the truth: the called allele contains the
trimmed, pre-mutation germline that was simulated, so that no method could
separate the two from the read. A junction is correct when it equals the
true sequence from Cys104 to Trp/Phe118. Junction accuracy is computed over
the reads whose two anchor codons are templated in the read: 9,953, 8,130
and 4,667 of 10,000 reads on \emph{easy}, \emph{mid} and \emph{hard}, and
1,441 of 3,000 at 30\,\% SHM.

\paragraph{D calls.} The trade-off between finding short D remnants and
leaving D-free junctions without a D can be moved at will: a tool that calls
a D more readily finds more short remnants, but it also calls a D on more
junctions that have none. D accuracy is therefore always reported together
with the share of the D-free reads that are given a D (``D-free given a D'';
the spurious D rate of the main text), and by the length of the remnant.

\subsection{Tool settings}
\label{sec:protocol-tools}

\paragraph{IgBLAST 1.22.0.} On the rhesus panels, IgBLAST ran with
\code{-organism rhesus\_monkey}, \code{-ig\_seqtype Ig} and
\code{-outfmt 19}, with the GenAIRR J auxiliary file and a rhesus IGHC
region database. On the human GenAIRR sets it ran with
\code{-organism human}, \code{-ig\_seqtype Ig} (\code{TCR} for TRB) and
\code{-outfmt 19}, on \code{makeblastdb} databases of the same FASTA files
and with a J auxiliary file. IGK and IGL have no D, but \code{igblastn}
requires a D database; it was given the IGH D database, which it does not
use for these loci. On the IgSim reads, IgBLAST ran with
\code{-organism human} on custom databases of the IMGT set. Its auxiliary
file lacked IGHJ5*03 and IGHJ5*04; they were added with the frame and CDR3
end of IGHJ5*01, which match the anchors of \swig.

\paragraph{MiXCR 4.7.0.} MiXCR ran as IgBench runs it: \code{buildLibrary}
on the GenAIRR FASTA files, then \code{align} with the
\code{generic-amplicon} preset and floating alignment boundaries, then
\code{exportAirr}. The left boundary floats at the start of the V. The right
boundary floats at the C or, on the human reads, which end in the J, at the
end of the J. MiXCR requires a licence from MiLaboratories
(\code{MI\_LICENSE} or a \code{mi.license} file). On the IgSim reads, MiXCR
ran with the options of the human GenAIRR sets, on a \code{buildLibrary}
library of the same IMGT FASTA. \code{buildLibrary} drops every sequence
that contains a wildcard base. Six V alleles of the IMGT set contain 1--4
unsequenced N bases, and they are the true V of 1.3--1.5\,\% of the reads.
So that MiXCR could call every allele, these bases were filled in MiXCR's
copy from the closest allele. On IgSim, MiXCR is scored by its AIRR call,
because germline-equivalent scoring already credits an allele that the read
cannot tell apart from the truth.

\paragraph{\swig{} 0.39.7.} \swig{} ran with its default calling profile
(R-optimized). On the rhesus panels it started from a prepared reference,
and on the human GenAIRR sets from FASTA files that give the anchor of each
allele in its header. On the IgSim reads it started from plain FASTA.

\paragraph{D acceptance rules of the other tools.} Each of the sequential
tools accepts a D by a fixed rule. IgBLAST 1.22 by default requires at least
5 consecutive matching D bases. MiXCR 4.7.0 drops D alignments that score
below 25, the score of five matched bases, and keeps further D hits only if
they score within 85\,\% of the best.
IMGT/JunctionAnalysis~\cite{yousfimonod2004junction} allows at most a set
number of D mismatches, by default 4 for IGH and 0 for TRB.

\subsection{Statistics}
\label{sec:protocol-stats}

Two tools, or two settings of \swig, are always compared on identical reads.
The command \code{swig evaluate} scores any table in the AIRR Rearrangement
format and therefore serves every tool alike. For each metric it computes
the difference in accuracy with its 95\,\% paired Wald interval, and an
exact two-sided McNemar test~\cite{mcnemar1947} on the reads that only one
of the two gets right. We report the differences and the McNemar
$p$-values. Posteriors are scored by reliability bins, the expected
calibration error~\cite{naeini2015ece}
$\text{ECE}=\sum_b n_b\,|\bar p_b-\bar e_b|/n$, where bin $b$ holds $n_b$ of
the $n$ reads, $\bar p_b$ is their mean reported probability and
$\bar e_b$ the observed frequency of the event, and the Brier
score~\cite{brier1950}.

\subsection{Hardware and timing}
\label{sec:protocol-timing}

Every run, for accuracy and for speed, used one machine: a laptop with an
Intel Core i9-13980HX processor, which has eight performance cores with two
hardware threads each and 16 efficiency cores and supports AVX2 but not
AVX-512. The machine had 62\,GB of memory and ran Linux on mains power.
Timed runs were made one at a time. Each tool's own thread option was set to
1, 4 or 8, and every command was restricted to the performance cores, so
that the three tools ran on the same cores. On the rhesus panels and the
human GenAIRR sets, the tools shared the 16 hardware threads of the eight
performance cores; for the IgSim timings of main-text \cref{M-tab:speed},
each thread had a performance core of its own. The rhesus panels were run on
eight threads and the human GenAIRR sets on four.

Unless stated otherwise, a timing is the wall-clock time of the whole
command: \code{igblastn} for IgBLAST; \code{align}, including the start-up
of Java, for MiXCR; and the command-line program, including the preparation
of the reference, for \swig. This preparation takes 25--28\,ms for the IMGT
IGH set, with every allele annotated on load. Each command also has a fixed
cost: a one-read run took 0.05\,s for \swig, 0.4--1.5\,s for IgBLAST
(0.05\,s with its usage report switched off) and 1.7--4.6\,s for MiXCR
across the timing passes.

The machine ran a desktop session, and the timed commands had network
access, so their times include IgBLAST's usage report to NCBI and MiXCR's
online licence check. The IgSim timings of main-text \cref{M-tab:speed} are
therefore provisional. A second pass on \emph{hard}, with the performance
cores shared, gave \swig{} and IgBLAST times within 7\,\% of the reported
ones on four and eight threads, and it gave MiXCR, whose Java runtime then
had spare cores, times that were 16--24\,\% shorter. The other MiXCR and
IgBLAST runs of that pass waited on network timeouts and were discarded; on
a slow connection, MiXCR's network calls held some runs at 74\,s whatever
their thread count. Other hardware may rank the thread scaling of the tools
differently.

% ===========================================================================
\section{Further results}
\label{sec:supresults}

These results test the model's own terms. Except for \cref{sec:dpriors},
they use the IgSim sets, the IMGT human IGH germline and the paired
statistics of \cref{sec:protocol}.

\subsection{D-prior design on human GenAIRR reads}
\label{sec:dpriors}

The D trim and N/P-length priors were designed on the IGH sets of the human
GenAIRR grid (\cref{sec:protocol-genairr}, main-text
\cref{M-tab:genairr-human}; GenAIRR 2.3.1~\cite{konstantinovsky2024genairr}
with its OGRDB reference~\cite{lees2020ogrdb}), not on the IgBench rhesus
panels, and calls are scored by the same Jaccard index against GenAIRR's tie
sets. The previous release's D-error sample and the prototype measurements
below are from that release's benchmark, and the instruction count from the
0.39.7 release measurements (callgrind on 1,000 IGH reads).

\paragraph{Where the D gain comes from.} Without its D trim and N/P-length
priors \swig{} 0.39.7 makes the D calls of its previous release, which was
level with IgBLAST at 2 and 5\,\% SHM on this grid (\cref{fig:dpriors}). An analysis of that release's errors on a
2,000-read sample of the 5\,\% SHM set found that 95 of its 103 wrong-gene D
calls overlapped the true D: the D had been found, but a short exact match of
another gene inside a mutated true D (for example \code{CTACAAT} inside
\code{ACTACAATAA}) outscored it by local alignment. A post-processing
prototype that held the placement graph's V end and J start fixed and
re-decided the D under a posterior over allele, offset and both trims showed
that the gain came from the trim and N-length priors and not from a
different decision rule or hypothesis space: with flat priors it matched the
graph's D calls (93.12 against 93.16 at 5\,\% SHM), and restricting each
allele to one offset, as the graph does, cost at most 0.1 point. Moving both
priors into the graph as cut terms (\cref{sec:objective}) recovered most of
the prototype's gain (95.42 against 95.67 at 5\,\% SHM) at 0.6\,\% more
instructions, whereas the prototype as a separate pass would have added
37--47\,\% to the run time.

\begin{figure}[t]
\centering
\begin{tikzpicture}
\begin{axis}[width=0.78\linewidth,height=5.6cm,
  xlabel={SHM rate (\%)},ylabel={Human IGH D score (\%)},
  xtick={0,2,5,10},ymin=80,ymax=100,xmin=-0.5,xmax=10.5,
  legend style={font=\scriptsize, at={(0.98,0.98)}, anchor=north east,
    draw=none, fill=white, fill opacity=0.9, text opacity=1},
  tick label style={font=\scriptsize},label style={font=\small},
  grid=major,grid style={gray!20},mark size=1.8pt,line width=0.9pt]
\addplot[cSwig,mark=*] coordinates {(0,98.47) (2,97.28) (5,95.42) (10,90.24)};
\addlegendentry{\swig{} 0.39.7}
\addplot[cOff,mark=square*,dashed] coordinates {(0,98.17) (2,96.38) (5,93.16) (10,87.12)};
\addlegendentry{\swig, D priors off}
\addplot[cIgb,mark=triangle*] coordinates {(0,97.80) (2,96.55) (5,93.24) (10,84.28)};
\addlegendentry{IgBLAST 1.22}
\end{axis}
\end{tikzpicture}
\caption{Design check: human IGH D accuracy on the GenAIRR SHM grid with
and without the D priors (mean Jaccard score; 10,000 reads per rate).
``D priors off'' is \swig{} 0.39.7 with \code{-{}-d-trim off -{}-np-length
off}. MiXCR 4.7.0 on the same reads is in main-text
\cref{M-fig:genairr-human,M-tab:genairr-human}.}
\label{fig:dpriors}
\end{figure}

Because the priors were designed on that human grid, we simulated four IGH
scenarios again from a disjoint part of GenAIRR's random stream, sharing no
read with the design sets (\cref{tab:oos}). The picture is unchanged: the
priors add 0.5--3.4 points of D, at least as much as on the design sets, and
\swig{} leads IgBLAST and MiXCR 4.7.0 on V, D and J in every scenario, on D
by 0.9--6.3 points over IgBLAST and 2.0--7.7 over MiXCR's tied hits (its
AIRR call: 75.0--87.8\,\% D, 94.9--99.6\,\% V).

\begin{table}[tbp]
\centering
\caption{Out-of-sample check for the D priors: GenAIRR human IGH reads
disjoint from the design grid of \cref{fig:dpriors} (10,000 per scenario;
mean Jaccard score, \%). ``Priors off'': \swig{} 0.39.7 with
\code{-{}-d-trim off -{}-np-length off}. MiXCR 4.7.0 counts every hit tied
with its top score.}
\label{tab:oos}
\footnotesize
\setlength{\tabcolsep}{3.2pt}
\begin{tabular}{@{}l rrr r rrr rrr@{}}
\toprule
 & \multicolumn{3}{c}{\swig{} 0.39.7} & priors off & \multicolumn{3}{c}{IgBLAST 1.22} & \multicolumn{3}{c}{MiXCR 4.7.0} \\
 & & & & & & & & \multicolumn{3}{c}{tied hits} \\
\cmidrule(lr){2-4}\cmidrule(lr){5-5}\cmidrule(lr){6-8}\cmidrule(l){9-11}
Scenario & V & D & J & D & V & D & J & V & D & J \\
\midrule
no SHM     & 99.98 & 98.40 & 99.65 & 97.89 & 99.82 & 97.54 & 99.48 & 99.74 & 96.43 & 99.39 \\
5\,\% SHM  & 99.66 & 94.97 & 99.39 & 92.71 & 99.59 & 92.65 & 99.15 & 99.42 & 91.00 & 99.03 \\
10\,\% SHM & 99.46 & 90.20 & 98.97 & 86.82 & 99.44 & 83.94 & 98.47 & 98.88 & 82.46 & 98.07 \\
realistic  & 99.54 & 94.75 & 99.33 & 92.43 & 99.44 & 92.12 & 99.13 & 98.85 & 90.48 & 98.86 \\
\bottomrule
\end{tabular}
\end{table}

\subsection{Truncated reads and the open-flank cost}
\label{sec:truncated}

Real reads can stop at their D (short or 3$'$-truncated reads) or start
inside CDR3 (5$'$-truncated reads), and a random tail beyond a V or before a J offers a D search
window where no D exists. \Cref{tab:truncated} cuts IgSim reads accordingly.
With the open-flank cost, a D at the read's open end is called right on 88\,\%
(\emph{mid}) and 82\,\% (\emph{hard}) of reads cut at the D's end and on 94 and
83\,\% of reads that start after the V, while a V followed by 100 random bases
receives a D on 2.6 and 11.2\,\% of reads, against 88--93\,\% without either
mechanism. Plausibility constraints alone, without the cost, remove the
chance Ds but also every real D at an open end. Chance Ds within about 30
bases of the V remain (68 and 77\,\%): they sit where a real D of a truncated
read sits, and no length or window term can separate them.

\begin{table}[tbp]
\centering
\caption{Truncated IgSim reads (2,000 per set and preset): right D calls on
reads that hold a D, and any D call on D-free constructs. ``Cost'' is the
open-flank cost, ``guards'' the plausibility constraints; the default has both.}
\label{tab:truncated}
\footnotesize
\begin{tabular}{@{}l l l rrrr@{}}
\toprule
Set & Read & Measure & default & cost off & guards off & both off \\
\midrule
\multicolumn{7}{@{}l}{\emph{mid}}\\
dcut      & V\,N\,D, cut at the D's end   & right D (\%) & 88.0 & 0.0 & 86.5 & 86.5 \\
cdr3start & N\,D\,N\,J\ldots, V removed     & right D (\%) & 93.6 & 0.0 & 93.5 & 93.5 \\
vcut30    & V + 30 random bases           & any D (\%)   & 67.5 & 0.1 & 66.0 & 66.0 \\
vcut100   & V + 100 random bases          & any D (\%)   & 2.6  & 0.2 & 5.7  & 88.3 \\
jcut100   & 100 random bases + J\ldots    & any D (\%)   & 9.7  & 0.1 & 10.4 & 91.8 \\
\addlinespace
\multicolumn{7}{@{}l}{\emph{hard}}\\
dcut      & V\,N\,D, cut at the D's end   & right D (\%) & 81.5 & 0.0 & 79.3 & 79.3 \\
cdr3start & N\,D\,N\,J\ldots, V removed     & right D (\%) & 83.2 & 0.0 & 83.2 & 83.2 \\
vcut30    & V + 30 random bases           & any D (\%)   & 76.6 & 0.2 & 74.4 & 74.5 \\
vcut100   & V + 100 random bases          & any D (\%)   & 11.2 & 0.8 & 14.8 & 93.2 \\
jcut100   & 100 random bases + J\ldots    & any D (\%)   & 27.2 & 0.0 & 28.2 & 90.6 \\
\bottomrule
\end{tabular}
\end{table}

\subsection{What each modelling term contributes}
\label{sec:ablation}

\Cref{tab:ablation} switches each term off on the IgSim \emph{mid} and
\emph{hard} sets and compares with the default read by read. The D trim and
N/P priors raise D accuracy on both presets ($+0.37$ and $+1.08$ points;
$p=0.02$ and $3\cdot10^{-9}$) and, on \emph{mid}, also leave 6.1 points more
D-free reads without a D. The fixed $(2,-3)$ scale (\code{-{}-calls fixed})
cannot carry the D trim and N/P terms, so its call changes the scale and
drops those priors together; its posterior still reads the evidence at the
read's divergence with the priors, but with presence costs constant in nats. Against the run
with the priors off, whose call differs only in the call model, calling at
each read's own divergence leaves 3.0 points more D-free reads without a D
on \emph{mid} (24.5 against 27.6\,\%, $p=0.002$; not significant on
\emph{hard}), but gets 0.2--0.3 points fewer D alleles right (93.02 against
93.24\,\% on \emph{mid}, $p=0.06$; 85.83 against 86.09\,\% on \emph{hard},
$p=0.014$) and 0.2--0.3 points fewer junctions exact (99.51 against
99.78\,\% on \emph{mid}, $p=5\cdot10^{-7}$; 96.87 against 97.09\,\% on
\emph{hard}, $p=0.03$). By itself it buys specificity at a small cost in
accuracy; its value lies in making the call a mode of the reported
posterior and in the D priors it makes possible. The priors more than
recover the D-allele loss on \emph{hard} (0.82 points above the fixed-scale
call; not significant on \emph{mid}), but not the junction loss: their effect on junctions is not significant
(99.53 against 99.51\,\% on \emph{mid}, 96.81 against 96.87\,\% on
\emph{hard}; \cref{tab:ablation}), and the fixed-scale call keeps 0.25 and
0.28 points more junctions exact than the default. The open-flank cost and
the conserved-anchor D window change essentially nothing on full-length
reads: they act on reads that stop in or start after the junction, which the
\emph{mid} and \emph{hard} presets rarely produce (\cref{tab:truncated} for
the cost). The max-marginal D decoder changes the D allele of a few dozen
full-length reads but gains exactly as many as it loses on IgSim (6 against
6 on \emph{mid}, 26 against 26 on \emph{hard}); on the human GenAIRR grid
(\cref{sec:dpriors}) it adds up to 0.2 points of IGH D (95.55 against 95.42 at 5\,\% SHM), too little to make
it the default. Skipping posteriors leaves every call unchanged. Raising the
prior odds of a D (\code{sensitive-d}, prior odds 2) gains 0.8--1.0 points of
D for 6--17 points more D-free reads given a D (\cref{M-fig:operating}). The
profiles without a D prior give a D to 64--91\,\% of D-free reads. On
\emph{hard} both are dominated and lose 0.8--1.9 points of D; on \emph{mid}
\code{sensitive-d} dominates the \code{truth} profile, and the
\code{igblast} profile's 0.1-point edge in D over \code{sensitive-d} (94.55
against 94.43\,\%) costs 51 points more D-free reads given a D.

\begin{table}[tbp]
\centering
\caption{Ablations on IgSim (10,000 reads per preset). $\Delta$: change
against the default in percentage points (positive: better), paired on the
same reads; ``0'': no read changes correctness; ``n.s.'': $p\ge0.05$ (exact
McNemar). ``D-free given a D'' and
the D-presence Brier score are absolute (default: 18.4\,\% and 0.0217 on
\emph{mid}, 15.8\,\% and 0.0227 on \emph{hard}). The fixed-scale row's
posterior still reads the evidence at the read's divergence with the D
priors, but charges each presence cost as a constant number of nats rather
than as $P^{*}_k/m$ matched columns, so its Brier score differs slightly from
the default's.}
\label{tab:ablation}
\footnotesize
\setlength{\tabcolsep}{3.5pt}
\begin{tabular}{@{}l rrrr rrrr@{}}
\toprule
 & \multicolumn{4}{c}{\emph{mid}} & \multicolumn{4}{c}{\emph{hard}} \\
\cmidrule(lr){2-5}\cmidrule(l){6-9}
Change & $\Delta$D & D-free & $\Delta$junction & Brier & $\Delta$D & D-free & $\Delta$junction & Brier \\
 &  & given a D & & D pres. & & given a D & & D pres. \\
\midrule
D trim and N/P priors off       & $-0.37$ & 24.5 & n.s.     & 0.0240 & $-1.08$ & 16.8 & n.s.     & 0.0218 \\
fixed-scale call (no D priors)  & n.s.    & 27.6 & $+0.25$  & 0.0224 & $-0.82$ & 16.4 & $+0.28$  & 0.0222 \\
open-flank cost off             & 0       & 18.4 & 0        & 0.0217 & $-0.06$ & 15.8 & 0        & 0.0239 \\
plausibility constraints off    & 0       & 18.4 & $+0.44$  & 0.0217 & n.s.    & 15.8 & $+2.85$  & 0.0223 \\
conserved-anchor window off     & n.s.    & 18.4 & n.s.     & 0.0216 & n.s.    & 16.0 & 0        & 0.0222 \\
max-marginal D allele           & n.s.    & 18.4 & n.s.     & 0.0217 & n.s.    & 15.8 & 0        & 0.0227 \\
no posteriors                   & 0       & 18.4 & 0        & --     & 0       & 15.8 & 0        & --     \\
\addlinespace
profile \code{sensitive-d}      & $+1.03$ & 35.0 & 0        & 0.0232 & $+0.84$ & 22.1 & n.s.     & 0.0217 \\
profile \code{truth}            & n.s.    & 64.5 & $-0.39$  & 0.0423 & $-1.87$ & 69.4 & $-1.52$  & 0.0417 \\
profile \code{igblast}          & $+1.16$ & 86.1 & $-0.30$  & 0.0474 & $-0.78$ & 91.1 & $-2.21$  & 0.0458 \\
\bottomrule
\end{tabular}
\end{table}

\subsection{Where the remaining errors are}
\label{sec:errors}

\Cref{tab:derrors} attributes every wrong or missing D call on 10,000 IgSim
reads per preset to the step where it arose. With a remnant of at least 8
nucleotides, the true D is almost never outside the search windows (0.04\,\%
on \emph{hard}) and almost never covered by the called J or V, the failure
mode of sequential annotation. Most errors are adjudication errors: a
placement of a correct allele existed, but the path chose another allele.
These are short, mutated remnants that another allele explains as well or
better under the model. On remnants of 1--7 nucleotides about half of the
errors (53 and 58\,\%) are a missing D: in 38 and 40\,\% of the errors
adjudication declined a D although a placement of an equivalent allele
existed, as the D prior intends for evidence this weak, and in 16 and
17\,\% no such placement survived the D search.

The same holds for V and J. Aligning every allele unbanded under the same
objective (an exhaustive oracle) found no V search error in 2,000 reads per
preset (from the same sample as \cref{tab:derrors}) and 1 and 3 J search
errors on \emph{mid} and \emph{hard}; the
remaining V errors are ties (26 and 84 reads) and model errors (12 and 68),
in which the exhaustive optimum itself is not the truth because mutations
fall on the positions that separate alleles.

\begin{table}[tbp]
\centering
\caption{Attribution of D errors (\% of reads with a D remnant in the class),
on a separate sample of 10,000 IgSim reads per preset (\cref{M-tab:igsim,tab:ablation,tab:calibration}
use another sample). ``Adjudication'': a placement of an
equivalent allele existed but the path took another D or none; ``not
found'': no such placement survived the D search.}
\label{tab:derrors}
\footnotesize
\begin{tabular}{@{}l rr rr@{}}
\toprule
 & \multicolumn{2}{c}{remnant $\ge 8$\,nt} & \multicolumn{2}{c}{remnant 1--7\,nt} \\
\cmidrule(lr){2-3}\cmidrule(l){4-5}
 & \emph{mid} & \emph{hard} & \emph{mid} & \emph{hard} \\
\midrule
reads                                      & 8,535 & 8,483 & 974 & 1,006 \\
D correct                                  & 97.90 & 91.76 & 51.85 & 44.83 \\
wrong D, adjudication                      & 1.91 & 6.90 & 17.25 & 16.60 \\
wrong D, not found                         & 0.12 & 0.83 & 5.34 & 6.76 \\
wrong D, outside the windows               & 0.00 & 0.04 & 0.00 & 0.00 \\
missing D, adjudication                    & 0.04 & 0.22 & 18.07 & 22.27 \\
missing D, not found                       & 0.04 & 0.26 & 7.49 & 9.54 \\
called J covers $\ge 3$ bases of the true D & 0.04 & 0.14 & 0.51 & 0.99 \\
\bottomrule
\end{tabular}
\end{table}

\subsection{Calibration}
\label{sec:calibration}

Posteriors are well calibrated wherever the candidates contain the truth
(\cref{fig:calibration}). The expected calibration error is at most 0.5\,\%
for V and 0.8\,\% for J on the \emph{easy}, \emph{mid} and \emph{hard}
presets, and 1.9\,\% for the D allele and
1.1\,\% for D presence on \emph{hard} (\cref{tab:calibration}). The D-allele
posterior is over-confident in the 0.5--0.95 range: reported probabilities
of 85\,\% correspond to 77--78\,\% observed. At 30\,\% SHM the V and J
posteriors are over-confident (ECE 9.8 and 13.2\,\%): when retrieval misses
the true gene, no candidate can carry its mass.

\begin{table}[tbp]
\centering
\caption{Calibration of the reported posteriors on IgSim reads: expected
calibration error (\%) and Brier score.}
\label{tab:calibration}
\footnotesize
\begin{tabular}{@{}l rrrr rrrr@{}}
\toprule
 & \multicolumn{4}{c}{ECE (\%)} & \multicolumn{4}{c}{Brier} \\
\cmidrule(lr){2-5}\cmidrule(l){6-9}
Preset & V & D & J & D presence & V & D & J & D presence \\
\midrule
easy       & 0.09 & 0.20 & 0.05 & 0.04 & 0.0004 & 0.0010 & 0.0006 & 0.0001 \\
mid        & 0.12 & 0.96 & 0.08 & 0.76 & 0.0064 & 0.0292 & 0.0063 & 0.0217 \\
hard       & 0.47 & 1.87 & 0.80 & 1.14 & 0.0278 & 0.0556 & 0.0343 & 0.0227 \\
30\,\% SHM & 9.78 & 20.91 & 13.22 & 2.53 & 0.1654 & 0.2340 & 0.1978 & 0.0488 \\
\bottomrule
\end{tabular}
\end{table}

\begin{figure}[t]
\centering
\begin{tikzpicture}
\begin{groupplot}[group style={group size=2 by 1,horizontal sep=1.4cm},
  width=0.46\linewidth,height=5.4cm,xmin=0,xmax=100,ymin=0,ymax=100,
  xlabel={reported posterior (\%)},
  tick label style={font=\scriptsize},label style={font=\small},title style={font=\small},
  grid=major,grid style={gray!20},every axis plot/.append style={mark size=1.6pt,line width=0.8pt},
  legend style={font=\scriptsize, at={(0.97,0.03)}, anchor=south east,
    draw=none, fill=white, fill opacity=0.9, text opacity=1}]
\nextgroupplot[title={\emph{mid}},ylabel={observed frequency (\%)}]
\addplot[gray,dashed,thin,forget plot] coordinates {(0,0) (100,100)};
\addplot[cV,mark=*] coordinates {(86.91,86.36) (93.56,91.61) (97.80,97.57) (99.69,99.78) (99.98,99.98)};
\addlegendentry{V}
\addplot[cD,mark=square*] coordinates {(39.24,47.93) (66.65,60.40) (85.08,76.54) (92.79,87.88) (97.87,97.42) (99.58,99.66) (99.98,99.90)};
\addlegendentry{D}
\addplot[cJ,mark=triangle*] coordinates {(72.72,80.77) (85.55,91.30) (93.44,94.12) (97.90,97.83) (99.67,99.60) (99.98,99.96)};
\addlegendentry{J}
\addplot[black,mark=o] coordinates {(16.46,23.99) (66.83,61.45) (85.30,81.99) (92.74,86.62) (97.38,95.45) (99.62,99.06) (100.00,99.99)};
\addlegendentry{D presence}
\nextgroupplot[title={\emph{hard}}]
\addplot[gray,dashed,thin,forget plot] coordinates {(0,0) (100,100)};
\addplot[cV,mark=*] coordinates {(73.61,82.98) (86.83,88.10) (93.13,93.88) (97.20,97.01) (99.59,100.00) (99.99,99.66)};
\addplot[cD,mark=square*] coordinates {(39.01,45.17) (65.79,57.35) (85.45,78.07) (93.08,90.41) (97.31,97.01) (99.62,98.37) (99.98,99.55)};
\addplot[cJ,mark=triangle*] coordinates {(70.35,71.07) (86.74,85.08) (93.29,92.72) (97.28,96.57) (99.57,99.19) (99.99,98.91)};
\addplot[black,mark=o] coordinates {(17.92,32.44) (66.34,73.65) (85.90,88.19) (92.80,90.18) (97.53,96.62) (99.59,99.85) (100.00,100.00)};
\end{groupplot}
\end{tikzpicture}
\caption{Reliability of the reported posteriors on IgSim reads: mean
reported probability against observed frequency of the event, per bin
(bins with at least 20 reads; most reads fall in the top bins). The event of
a segment is ``the reported names include the read's allele up to germline
equivalence'', that of D presence ``the read has a simulated D''.}
\label{fig:calibration}
\end{figure}

\subsection{Constant region}
\label{sec:constant-results}

On 10,000 IgSim reads per preset (both strands), about 80\,\% of which
continue into one of the 101 human IGHC alleles with a complete coding
start, the primary C call
is anchored-equivalent to the truth on 98.5--99.2\,\% of reads with 16--31
nucleotides of C and on at least 99.65\,\% from 32 nucleotides, and C was
called on at most 0.05\,\% of the about 2,000 reads per preset that stop
before C. C parts of 8--15 nucleotides are called on 53--60\,\% of reads, as
the 12-column floor intends. V, D and J calls are identical with and
without the C stage by construction.

\subsection{Throughput on the human GenAIRR sets}
\label{sec:genairr-speed}

\Cref{tab:genairr-speed} gives the throughput of the three tools on the
human GenAIRR sets of main-text \cref{M-tab:genairr-human}, each set run once
through the accuracy harness with its options on four threads, the tools
sharing the 16 hardware threads of the performance cores. \swig{} assigned
15,400--28,500 reads per second, 86--153 times IgBLAST's throughput and
8.3--13.4 times MiXCR's. The timed runs reproduced the scored outputs byte
for byte. The IGH 10\,\% SHM and realistic sets ran while a compile used
other cores and took up to 36\,\% longer than in the two accuracy runs of the
same commands, which were not restricted to the performance cores, so they
are left out. On the other sets the timed \swig{} and MiXCR runs, and
IgBLAST's on IGH and TRB, were within 7\,\% of the accuracy runs, except
MiXCR on IGH 5\,\% SHM (21\,\% longer; marked).

\begin{table}[tbp]
\centering
\caption{Throughput on the human GenAIRR sets (10,000 reads each; reads per
second; wall clock of the whole command, one run per set; four threads).
\swig's includes database preparation and MiXCR's its Java start-up; MiXCR's
\code{buildLibrary} step is outside the timing. The last two columns are
\swig's throughput relative to each tool's. $^\dagger$21\,\% slower than in
the accuracy runs of the same command (1,690--1,720 reads per second).}
\label{tab:genairr-speed}
\footnotesize
\begin{tabular}{@{}l rrr rr@{}}
\toprule
Set & \swig{} & IgBLAST 1.22 & MiXCR 4.7.0 & vs IgBLAST & vs MiXCR \\
\midrule
IGH, no SHM      & 15,379 & 179 & 1,480 & 86  & 10.4 \\
IGH, 2\,\% SHM   & 15,737 & 182 & 1,553 & 87  & 10.1 \\
IGH, 5\,\% SHM   & 15,813 & 184 & 1,423$^\dagger$ & 86  & 11.1 \\
IGK, 5\,\% SHM   & 27,924 & 194 & 2,076 & 144 & 13.4 \\
IGK, realistic   & 28,303 & 198 & 2,306 & 143 & 12.3 \\
IGL, 5\,\% SHM   & 27,810 & 181 & 2,088 & 153 & 13.3 \\
IGL, realistic   & 28,491 & 189 & 2,148 & 151 & 13.3 \\
TRB, no SHM      & 24,983 & 233 & 2,909 & 107 & 8.6  \\
TRB, realistic   & 24,975 & 243 & 3,013 & 103 & 8.3  \\
\bottomrule
\end{tabular}
\end{table}

\subsection{Kernel backends}
\label{sec:kernels}

\Cref{tab:kernels} times the alignment backends on one thread, on the
10,000-read IgSim sets and the machine of main-text \cref{M-tab:speed}. The
backends were timed in turn, three rounds each, so that slower periods of
the machine affect them alike. The default AVX2 kernel gives 2.3--2.4 times
the throughput of the scalar kernel, and the portable lanes, which leave
vectorisation to the compiler, 1.8--1.9 times. This CPU has no AVX-512, so
the AVX-512BW backend was not timed. Every backend and every thread count
produced byte-identical output on both sets.

\begin{table}[tbp]
\centering
\caption{Provisional one-thread throughput by alignment backend on 10,000
IgSim reads (reads per second; wall clock of the whole command including
database preparation; median of three runs, backends interleaved, on the
shared performance cores). A compile was running on other cores when the
pass began, so the ratios between rows are firmer than the rates; the AVX2
rates are within 5\,\% of the one-thread \swig{} rows of main-text
\cref{M-tab:speed}, timed in another pass. AVX2 is the default backend on
this CPU.}
\label{tab:kernels}
\footnotesize
\begin{tabular}{@{}l rr@{}}
\toprule
Kernel (1 thread) & \emph{mid} & \emph{hard} \\
\midrule
AVX2 (16 lanes)           & 3,010  & 3,266 \\
portable lanes (16)       & 2,401  & 2,483 \\
scalar kernel             & 1,302  & 1,335 \\
\bottomrule
\end{tabular}
\end{table}

% ===========================================================================
\FloatBarrier
\section{Limitations}
\label{sec:limitations}

This section expands the limitations that the Discussion of the main text
summarises (\cref{M-sec:discussion}).

\paragraph{Simulated reads.} All results are on simulated reads. Both
simulators make simplifying assumptions: allele usage is uniform, trims are
independent, and there are no chimeras or long adapters; IgSim also has no
hotspot-targeted SHM and no codon-length indels. Several design parameters
were measured on simulated reads: the open-flank allowance from simulated
junction lengths, the credit $\kappa$ on IgSim reads, and the D trim and N/P
priors from an analysis of the errors of the previous release on the human
GenAIRR grid (\cref{fig:dpriors}). The means of these priors (5, 7 and 8
bases) are generic rather than fitted, and GenAIRR's own trim tables did no
better in the prototype. On human GenAIRR reads disjoint from the design
sets the priors gained as much (\cref{tab:oos}), but those reads come from
the same simulator. The rhesus comparison (main-text \cref{M-tab:genairr})
uses a different species and SHM schedule. Validation on real repertoires,
for example from donors with sequenced germline genotypes, is the most
important next step.

\paragraph{Comparisons.} We did not compare with partis or IMGT/V-QUEST.
MiXCR's tied hits are derived from the hit scores that it exports; they are
not an output mode of MiXCR itself, whose AIRR export names one allele per
segment. On IgSim, MiXCR's copy of six IMGT V alleles has its N bases filled
in, because its library builder rejects them (\cref{sec:protocol-tools}).
Throughput was measured on one laptop, under the conditions described in
\cref{sec:protocol-timing}.

\paragraph{Plausibility constraints.} The constraints of \cref{sec:guards}
are inherited thresholds, not terms of the likelihood, and the call under
them is the best admissible subset optimum rather than an exact maximum over
every admissible path. They cost 0.4--2.9 points of junction exactness on
simulated reads (\cref{tab:ablation}), in exchange for robustness on
truncated, chimeric and cross-locus reads that the simulators do not
produce. A model of truncation and template switching would price what they
forbid.

\paragraph{D calls on D-free junctions.} Of the IgSim reads without a D
remnant, 16--18\,\% still receive a D, and a chance D within about 30 bases
of a truncated read's V is indistinguishable from a real D there
(\cref{tab:truncated}). On TRB, whose two D genes are short and rarely
trimmed as much as the prior expects, the credit $\kappa$ gives D-free reads
a D more often.

\paragraph{Extreme mutation.} At 30\,\% SHM the strand test and the seed
retrieval lose intact $k$-mers: 4.4\,\% of the reads are left uncalled, most
of them (3.5\,\% of all reads) because the wrong strand was chosen, V
accuracy falls below IgBLAST's, and the V and J posteriors are
over-confident because the true gene is no longer among the candidates.

\paragraph{Statistics.} E-values use exact ungapped $\lambda$ and $K$, which
is anti-conservative for gapped scores, and the total database length, which
is conservative for near-identical alleles; they gate hits and are not
calibrated $p$-values. One divergence per read assumes uniform mutation
across V, D and J.

\paragraph{One alignment per D allele and window.} A D allele's optimal
local alignment in its window need not be its best clip at the chosen cuts
(\cref{tab:search}). This is the one remaining approximation of the
segmentation model.

\FloatBarrier
\bibliographystyle{unsrtnat}
\bibliography{references}

\end{document}
